\section{Twists, boundedness, and categorical duality}\label{sec:foundations}

The kernel calculations below take place in categories of all twisted
$D$-modules on bundle stacks.  Their two necessary foundations are a precise
choice of twist and a finiteness statement: at the irrational exponents in
question, restriction to a suitable quasi-compact open is an equivalence.
We establish these facts first, and then explain the resulting kernel
formalism, including its renormalized evaluation pairing.

\subsection{Root data, determinant lines, and level conventions}

All groups retain their given root data.  Fix dual pinnings and a theta
characteristic $\vartheta$ on $X$, together with an isomorphism
$\vartheta^{\otimes2}\simeq\Omega_X$.  For a connected reductive group $H$ write
\[
 Y_H=\operatorname{Bun}_H(X),\qquad
 L_H|_E=\det R\Gamma(X,\mathfrak h_E).
\]
Thus $\det H^1$ occurs with exponent $-1$.  Write $\rho_H$ for the half-sum
of positive \emph{coroots}, and put $P_0=\vartheta^{2\rho_H}$, regarded as an
$H$-bundle through its torus reduction.  This uses the integral cocharacter
$2\rho_H$ and makes sense for the actual group $H$.  In Grassmannian
constructions we rigidify the determinant line by replacing $L_H$ with
$L_H\otimes(L_H|_{P_0})^{-1}$.  A constant line does not change the sheaf of
twisted differential operators.  When tensoring by integral powers, either
rigidification gives the same formula after retaining the corresponding
constant one-dimensional factor.

For a line $L$ on a smooth stack $Y$ and $s\in\mathbb C$, write $D(L^s)$ for
the presentable DG category of twisted $D$-modules.  Here is the sign
convention.  On a smooth scheme, if a frame $e$ is replaced by $ae$, the
local differentiation operator changes by $-s\,d\log a$.  Equivalently,
on the frame torsor of $L$ we impose strong Kummer equivariance with the
connection
\begin{equation}\label{foundations:kummer}
 d-s\,d\log t,\qquad q=\exp(2\pi i s).
\end{equation}
The associated de Rham local system has monodromy $q$.  On stacks these
categories and their functors are defined by strong derived descent.
The opposite twist is $L^{-s}$.  Sums of twists use the tensor product of
frame torsors, or equivalently their ratio torsor; these conventions also
apply on singular auxiliary stacks by smooth descent and Kashiwara's lemma.

We use sheaf conventions for operations: $\omega_Y$ is the unit of
$\otimes^!$, $k_Y$ is the constant object, and a lower star denotes de Rham
direct image.  To express calculations with left $D$-modules on smooth
$Y$, write $M^{\langle Y\rangle}$ for the object corresponding to the left
module $M[\dim Y]$.  In particular,
\[
 \mathcal O_Y^{\langle Y\rangle}=\omega_Y.
\]
Inside these brackets, $!$-pullback is left crystalline pullback without
an additional shift.  The convention descends to smooth stacks, with their
stack dimension.  The exponential object of phase $\phi$ is the left
connection $d+d\phi$ in this dualizing normalization.  All tensor products,
functors, and mapping complexes are derived.  Categories are cocomplete
and $\mathbb C$-linear; functors between them are continuous unless an
adjoint or a partially defined functor is explicitly specified.

We now translate the shifted level of the problem into the four categories
used in the proof.  Use the basic invariant forms
$\operatorname{Kil}/(2h^\vee)$ and
$\operatorname{Kil}^{\vee}/(2\check h^\vee)$.  Their inverse forms give
long roots, respectively long coroots, squared length $2$.  Thus inverse
shifted forms have scalar parameters $a,d$ related by $d=1/(ra)$, where
$r$ is the lacing number.  Set
\begin{equation}\label{foundations:levels}
 a=-c,\qquad d=\frac1{ra}=-\frac1{rc},\qquad
 k_a=a-h^\vee,\qquad k_d=d-\check h^\vee.
\end{equation}
Our notation is
\begin{equation}\label{foundations:four-categories}
 \begin{aligned}
 A&=D\bigl(L_G^{-k_a/(2h^\vee)}\bigr),&
 A'&=D\bigl(L_G^{k_a/(2h^\vee)}\bigr),\\
 B'&=D\bigl(L_{G^\vee}^{-k_d/(2\check h^\vee)}\bigr),&
 B&=D\bigl(L_{G^\vee}^{k_d/(2\check h^\vee)}\bigr).
 \end{aligned}
\end{equation}
In the convention of \cite[Section~1.1.5]{BogdanovaBetti}, a shifted level
$c$ means the exponent $(c-h^\vee)/(2h^\vee)$.  Since
\[
 -\frac{k_a}{2h^\vee}
 =\frac{c-h^\vee}{2h^\vee}+1,
\]
tensoring by $L_G$ identifies the source category in that convention with
$A$.  Its target at shifted level $-1/(rc)$ is exactly $B$.
This is an integral translation by the determinant line; it requires
neither a square root of that line nor divisibility of the basic form in
the lattice of integral levels.  Because $c\notin\mathbb Q$, all four
exponents in \eqref{foundations:four-categories} are irrational.

We retain all connected components:
\begin{equation}\label{foundations:components}
 \pi_0(Y_H)=\pi_1(H)
 =X_*(T_H)/\mathbb Z\Phi_H^\vee.
\end{equation}
In particular, modification labels and dual weights always belong to the
lattices of the specified groups.  For semisimple $H$, its finite centre
acts trivially on $L_H$, since the adjoint representation has trivial
central action.  Central inertia therefore gives the usual decomposition
into character sectors.  These components and sectors are retained by the
constructions below.

\subsection{An irrational twist cuts off the unstable directions}

We prove the boundedness statement in a form that also applies to the
Levi subgroups occurring later.  Let $H$ be connected reductive and let
$V$ be a finite-dimensional self-dual representation.  Put
\[
 L(V)|_E=\det R\Gamma(X,V_E),\qquad
 Q(z,w)=\operatorname{Tr}_V(zw).
\]
Assume that $Q$ is positive definite on the real cocharacter space.
For a degree $u\in\pi_1(H)$ its rational image has a unique representative
$u_{\mathrm{cen}}\in X_*(Z(H)^0)\otimes\mathbb Q$.  We write $Q(z,u)$
for $Q(z,u_{\mathrm{cen}})$ when $z$ is central.

\begin{lemma}[Inertia vanishing]\label{foundations:inertia}
Let $Z$ be an algebraic stack carrying a line $L$.  Suppose a central
$\mathbb G_m$ in its inertia acts on $L$ with a fixed nonzero integral
weight $m$.  If $sm\notin\mathbb Z$, then $D_Z(L^s)=0$.
The assertion remains valid after restricting to a locally closed
substack and after taking a product with a spectator stack carrying any
external twist.
\end{lemma}

\begin{proof}
Pull an object to the frame torsor of $L$.  The weighted scaling of the
frame supplied by the inertial $\mathbb G_m$ is $2$-isomorphic to the
projection.  Strong Kummer equivariance consequently identifies the two
pullbacks of the object to the product with $\mathbb G_m$, one with a
constant factor and the other with the Kummer factor of exponent $sm$.
On a smooth chart, apply de Rham direct image along $\mathbb G_m$.
The nonintegral Kummer factor has zero de Rham cohomology.  For example,
in the Laurent-polynomial de Rham complex its differential has eigenvalues
$n-sm$ on $t^n$, all invertible.  The constant factor has the nonzero finite
complex $H^\bullet_{\mathrm{dR}}(\mathbb G_m)$, containing $\mathbb C$ as a
summand.  Tensoring with that complex detects zero, even for unbounded
complexes.  Hence the pullback object vanishes.  Smooth pullback and
frame-torsor pullback are conservative.  The calculation leaves all
spectator factors unchanged.
\end{proof}

\begin{proposition}[Bounded support in the categorical sense]\label{foundations:cutoff}
For $s\notin\mathbb Q$, the category $D_{Y_H}(L(V)^s)$ vanishes on every
component of nontorsion degree.  There is a quasi-compact open
$j:U\hookrightarrow Y_H$, contained in the union of the torsion-degree
components, such that
\begin{equation}\label{foundations:cutoff-equivalence}
 j^!:D_{Y_H}(L(V)^s)\xrightarrow{\ \sim\ }D_U(L(V)^s)
\end{equation}
with inverse $j_*$.  The same open works for $-s$, and the equivalence
persists after taking a product with any spectator stack and any external
twist there.
\end{proposition}

\begin{proof}
First suppose that $u$ has nonzero rational image.  Nondegeneracy on the
centre gives an integral central cocharacter $z$ with $Q(z,u)\ne0$.
Decompose $V$ into its central weight spaces.  Riemann--Roch gives the
weight of $z$ on $L(V)$ as
\begin{equation}\label{foundations:rr-weight}
 \sum_{\mu}\mu(z)\chi(X,V_{\mu,E})=Q(z,u).
\end{equation}
The term involving $1-g(X)$ is zero because $V$ is self-dual.
The right side is an integer, being the weight of an algebraic
$\mathbb G_m$-action on a line.  Lemma~\ref{foundations:inertia} proves the
first assertion, including its restriction to any locally closed bundle
condition in that component.

Now fix a torsion degree.  We use the characteristic-zero
Harder--Narasimhan theorem for principal bundles: canonical parabolic
reductions describe the strata, the HN types form a discrete set with
bounded denominators, and bounded HN regions admit quasi-compact open
neighbourhoods.  The associated bundle of a semistable Levi bundle under
a representation on which the connected centre acts by a character is
semistable.  These are the forms of the reduction and vanishing results
recalled in \cite[Sections~7.3--7.4 and~10.1]{DrinfeldGaitsgoryBun}.
Write $\nu$ for the dominant rational HN coweight.  Its central component
is zero because the degree is torsion.  Put
\[
 b_X=\max\{0,2g(X)-2\}.
\]
We claim that the twisted category on its stratum vanishes whenever
$\alpha_i(\nu)>b_X$ for some simple root $\alpha_i$.

Let $Q_i$ be the maximal standard parabolic obtained by omitting
$\alpha_i$, and let $M_i$ be its Levi subgroup.  The canonical parabolic
of $\nu$ is contained in $Q_i$.  Projecting the canonical reduction to
$M_i$ identifies the HN stratum with a stack of unipotent extensions over
the corresponding HN stratum of $Y_{M_i}$.  Filter the unipotent radical
of $Q_i$ by height in $\alpha_i$.  Each vector-group quotient, with its
filtration induced by the canonical reduction inside $M_i$, has
semistable associated graded pieces whose slopes are root evaluations
$\beta(\nu)$.  Every such root $\beta$ contains $\alpha_i$ with positive
coefficient; dominance therefore gives
\[
 \beta(\nu)\geq\alpha_i(\nu)>2g(X)-2.
\]
Serre duality and semistability imply $H^1=0$ for every piece, hence for
each successive extension bundle.  Riemann--Roch then makes the ranks of
$H^0$ constant on the fixed stratum.  Consequently these extension stacks
are, locally in families, successive classifying stacks of smooth
unipotent groups.  The section of split extensions is a smooth
surjective cover in the sense of stacks.

Choose a positive integral multiple $z$ of the $i$th fundamental
coweight, in the semisimple span, that belongs to $X_*(T_H)$.  It is
central in $M_i$.  On the split section, the determinant filtration
identifies $L(V)$ with the determinant line for the restricted
$M_i$-representation.  Formula~\eqref{foundations:rr-weight} now gives
its $z$-weight as $Q(z,\nu)$.  This is strictly positive: on each simple
factor the trace form is a positive multiple of the basic form, and the
inverse Cartan matrix has strictly positive entries.  The dominant
coweight $\nu$ is nonzero on the factor containing $\alpha_i$.
Lemma~\ref{foundations:inertia} thus kills the pullback to the split
section, and smooth descent kills the whole HN stratum.  The proof also
kills its restriction to any locally closed substack.

The remaining HN types satisfy $0\leq\alpha_i(\nu)\leq b_X$ for every
simple root.  Their central part is zero, so they lie in a bounded
region.  Bounded denominators make their number finite.  There are only
finitely many torsion degrees in $\pi_1(H)$.  Choose a quasi-compact HN
open $U$ containing all these types.  On the complement, restriction to
every HN stratum vanishes.  To pass from strata to arbitrary objects,
work on each bounded open in the HN exhaustion: its complement of $U$
has a finite HN stratification, and the localization triangles show
successively that its twisted category is zero.  The exhaustion detects
objects, so the same holds on the full complement.  The localization
triangle for $j$ then proves \eqref{foundations:cutoff-equivalence}, with
inverse $j_*$.  All arguments commute with an external spectator and
use only the nonintegrality of a nonzero integer times $s$; they apply
unchanged to $-s$.
\end{proof}

The assertion is about $!$-restriction of categories.  It makes no
set-theoretic claim about closures of supports of the underlying
complexes.  For the simple groups in the main theorem one takes
$V=\mathfrak h$.  For a Levi subgroup, the restriction of the original
adjoint representation supplies a self-dual representation with positive
trace form, including on the Levi centre.

\subsection{Operations with compatible twists}\label{foundations:operations}

A QCA stack is a quasi-compact algebraic stack with affine automorphism
groups; our stacks are locally of finite type over $\mathbb C$.
We use the $D$-module formalism of
\cite{DrinfeldGaitsgoryFiniteness}: compact generation, categorical
duality, and renormalized de Rham direct image with base change and the
projection formula.  A quasi-compact morphism is safe when the neutral
components of the automorphism groups in its geometric fibres are
unipotent.  Safe direct image agrees with ordinary de Rham $*$-image and
is continuous; schematic quasi-compact morphisms are examples
\cite[Sections~9--10]{DrinfeldGaitsgoryFiniteness}.

For correspondences with compatible twists the same formulas apply.
Compatibility identifies the relevant pullbacks of frame torsors, or
trivializes their ratio when twists are added.  Pulling to these torsors
reduces $!$-base change, the $!$-tensor projection formula, and continuity
of safe direct image to the ordinary formulas.  Their natural
identifications commute with frame scaling, so they descend with the
specified Kummer equivariance.  This also explains why one must retain
the compatibility of lines throughout a correspondence.

We use the six operations on locally holonomic objects wherever they
are defined.  In particular, for schematic separated finite-type maps,
shriek direct image on holonomic arguments is the partial left adjoint
to $!$-pullback, with the adjunction tested against arbitrary $D$-modules.
This assertion and Verdier duality, which negates the twist, can be
checked on smooth charts.  Riemann--Hilbert is used only for regular
holonomic objects, again on charts; Kummer connections are regular
singular.  For a smooth map $q$ of relative dimension $e$, our convention
is $q^*=q^![-2e]$.

\subsection{Compact generation and opposite-twist duality}

The cutoff has reduced our bundle categories to smooth QCA stacks.
We record explicitly why twisting preserves the finiteness and kernel
statements needed later.

\begin{proposition}\label{foundations:categories}
Let $Y$ be a smooth quasi-compact open in a bundle stack as above, and
let $\tau$ be a determinant-line twist.  Then $D_\tau(Y)$ is compactly
generated.  On every smooth finite-type scheme chart, the pullback of a
compact object has bounded coherent cohomology.  For two such
stacks and twists, external tensor product induces an equivalence
\begin{equation}\label{foundations:kunneth}
 D_{\tau_1}(Y_1)\otimes D_{\tau_2}(Y_2)
 \xrightarrow{\ \sim\ }
 D_{\tau_1\boxplus\tau_2}(Y_1\times Y_2).
\end{equation}
The ordinary categorical dual of $D_\tau(Y)$ is $D_{-\tau}(Y)$, with
evaluation and coevaluation
\begin{equation}\label{foundations:evaluation}
 (M,N)\longmapsto
 \Gamma_{\mathrm{dR}}^{\mathrm{ren}}(Y,M\otimes^!N),
 \qquad \Delta_*\omega_Y.
\end{equation}
The twist on the diagonal is canonically trivial.  These assertions
remain valid for the full bundle categories reduced by
Proposition~\ref{foundations:cutoff}.
\end{proposition}

\begin{proof}
\emph{The induction adjunction.}
We first specify the underlying quasi-coherent complex of a twisted
object.  Write $\tau=s\,d\log L$ and let $P=L^\times\to Y$ be the frame
torsor.  An object $N\in D_\tau(Y)$ is represented on $P$ by a strongly
Kummer-equivariant D-module.  Forgetting its connection in the normalized
left-module convention turns the Kummer connection
$d-s\,d\log t$ into the trivial line with its multiplicative frame $1$.
In this normalization, smooth $!$-pullback induces ordinary pullback of
the underlying quasi-coherent complexes, with no remaining dimension shift.
Its equivariance isomorphisms therefore give ordinary quasi-coherent
descent along $P\to Y$, including the descent coherences.  Denote the
result by $N^{\mathrm{oblv}}$.  This operation is conservative and
preserves colimits.  It retains the equivariance maps, rather than
replacing them by identity maps; in particular it retains their
infinitesimal isotropy actions.

If $Y$ has dimension $d$ and canonical line $\omega_Y^{\mathrm{can}}$,
this normalization is right-module oblivion tensored with
$(\omega_Y^{\mathrm{can}})^{-1}[-d]$ on line-trivializing charts.
We will construct $\operatorname{Ind}_\tau$ on $\operatorname{Perf}(Y)$
with the adjunction
\begin{equation}\label{foundations:induction-adjunction}
 \operatorname{RHom}_{D_\tau(Y)}
       (\operatorname{Ind}_\tau(F),N)
 \simeq R\Gamma(Y,F^\vee\otimes N^{\mathrm{oblv}}).
\end{equation}
Since perfect complexes compactly generate $\operatorname{QCoh}(Y)$,
this formula will give compact generators for $D_\tau(Y)$.  Its
external-product compatibility will give K\"unneth.  We verify the
evaluation pairing separately after these constructions.

\emph{Presentations and induction.}
The bounded Quot construction presents $Y$ as a quotient of a smooth
finite-type scheme by a product of general linear groups.  Indeed, a
faithful representation bounds the associated vector bundles; after a
uniformly large twist they have vanishing $H^1$ and are generated by
global sections.  Framing those sections gives a scheme atlas.  Reductions
of structure group are representable by schemes of sections because the
quotient by the reductive subgroup is affine.  We can therefore choose
a smooth presentation $p:Z\to Y$ which is a torsor for such a product.

Here is the twisted induction construction on this presentation.
It is the relative Lie-algebroid construction of
\cite[Section~6.3, especially Lemmas~6.3.15--6.3.16]{DrinfeldGaitsgoryFiniteness},
with the lift of the relative action to the twisting line retained.
For a smooth schematic chart $T\to Y$, let $A_T=T_{T/Y}$ be the relative
tangent Lie algebroid, with its anchor to $T_T$.  Relative differentiation
lifts $A_T$ to the Atiyah algebroid of the pulled-back line, and hence to
its twisted differential operators $\mathcal D_T^\tau$.  Explicitly, if
$a\in A_T$ has anchor $v_a$ and its lift acts on a local frame $e$ by
$b_ae$, the resulting operator is $v_a+s b_a$.  Under $e'=fe$ one has
$b'_a=b_a+v_a(\log f)$, while the differentiation generator changes by
$-s v_a(\log f)$; the operator is therefore independent of the frame.
For
$F\in\operatorname{Perf}(Y)$ define the chart expression
\begin{equation}\label{foundations:induction}
 \operatorname{Ind}_\tau(F)|_T
 =\left(\mathcal D_T^\tau
       \mathop{\otimes}^{L}_{U(A_T)}F_T\right)^{\langle T\rangle}.
\end{equation}
Here $F_T$ has the canonical relative connection coming from its
pullback from $Y$, and $U(A_T)$ denotes the enveloping algebra.  If the
anchor has a kernel, that kernel is retained: its action on the twisting
line supplies the corresponding scalar action in
$\mathcal D_T^\tau$.  Trivializing the line on $T$ does not discard this
isotropy action.

For a smooth map $T'\to T$ over $Y$, the relative tangent algebroid
$A_{T'}$ is the inverse-image algebroid.  Resolve first in
$T_{T'/T}$ using its relative Spencer complex.  The result is precisely
the transfer bimodule for smooth pullback applied to
\eqref{foundations:induction}.  Resolving successively for two smooth
maps gives the same comparison, by associativity of transfer tensor
products.  Thus the chart objects and their comparisons descend.
This argument remains valid when an anchor has a kernel: after smooth
base change by $Z\to Y$, it is the relative differentiation along
ordinary smooth morphisms over $Z$, with its descent data.

At zero twist, the induction used here is the induction of
\cite[Section~6.3.21]{DrinfeldGaitsgoryFiniteness}
applied, in its IndCoh normalization, to
$F\otimes\omega_Y^{\mathrm{can}}[d]$.
For example, on $BH$ one has $d=-\dim H$ and
$\omega_{BH}^{\mathrm{can}}=\det\mathfrak h$.
At zero twist these factors cancel the
$\det(\mathfrak h^\vee)[\dim H]$ in the atlas formula of
\cite[Section~6.4.8]{DrinfeldGaitsgoryFiniteness};
\eqref{foundations:induction} then gives precisely derived
$\mathfrak h$-coinvariants on the point atlas.  The relative-algebroid
construction uses the same normalization at nonzero twist.

\emph{The adjunction, including unbounded coefficients.}
The Spencer resolution and enveloping-algebra adjunction on charts give
the desired formula \eqref{foundations:induction-adjunction} locally.
We explain the descent assertion implicit here.  Restricting $N|_T$ to
$A_T$ gives the relative connection on the pullback of
$N^{\mathrm{oblv}}$, with its derived descent homotopies.  The twist is
canonically split along these relative derivatives.  The mapping
complex is consequently the totalization, over the smooth faces of the
\v Cech nerve of $p$, of the corresponding relative de Rham complexes
with coefficients in $F^\vee\otimes N^{\mathrm{oblv}}$.

This totalization works for unbounded coefficients as well.  Put
$E=p_*\Omega^\bullet_{Z/Y}$ and
$p_n:Z^{[n]}=Z^{\times_Y(n+1)}\to Y$.  The product formula for relative
forms and the affine projection formula give, for any coefficient
complex $M$,
\[
 Rp_{n,*}\bigl(\Omega^\bullet_{Z^{[n]}/Y}\otimes p_n^*M\bigr)
 \simeq E^{\otimes^L(n+1)}\otimes^L M.
\]
The cofaces insert the unit of $E$.
The invariant Reynolds functional on
regular functions of the group, extended by zero on positive-degree
forms, gives a chain map $E\to\mathcal O_Y$ splitting the unit.
Left and right invariance make it independent of torsor
trivializations.  Applying this retraction in the first tensor factor
contracts the augmented tensor-power diagram.  Tensoring this explicit
contraction with any coefficient complex proves the asserted
unbounded totalization formula: the homotopy lowers cosimplicial degree
by one, so it acts on the product totalization without an infinite sum.
It contracts the augmented diagram,
not the individual de Rham complex of the group.

The forgetful functor is conservative and preserves colimits.  Perfect
complexes compactly generate $\operatorname{QCoh}(Y)$ on these smooth
quotient stacks, and their external products generate on products
\cite[Sections~3--4]{DrinfeldGaitsgoryFiniteness}.  Formula
\eqref{foundations:induction-adjunction} therefore makes the induced
perfects compact generators of $D_\tau(Y)$.  Their finite Spencer
resolutions on charts consist of coherent $\mathcal D_T^\tau$-modules.
They therefore have bounded coherent cohomology on each smooth
finite-type chart.  Finite cones and retracts preserve this property,
so it holds for all compact objects.  Induction commutes with external
products.  Applying
\eqref{foundations:induction-adjunction} to external perfects, together
with the quasi-coherent K\"unneth formula, proves full faithfulness on
compact generators in \eqref{foundations:kunneth}; these generators also
generate the target.  Hence \eqref{foundations:kunneth} is an equivalence.

\emph{The evaluation pairing.}
For an untwisted QCA stack, the evaluation is renormalized de Rham
cohomology and the coevaluation is $\Delta_*\omega_Y$
\cite[Sections~8.4 and~9.1]{DrinfeldGaitsgoryFiniteness}.
For opposite twists use the same two formulas, with the ratio line
trivialized on the diagonal.  To check either triangle identity, pull
the output leg to a smooth chart on which the output twisting line is
trivialized.  Base change replaces the diagonal by the graph of the
chart, and the projection formula cancels the two twists along this
graph.  The remaining expression is the untwisted graph calculation.
The graph is schematic and safe, so its renormalized image is its
ordinary star image and the composite is the identity.  These
identifications descend, since they use only the canonical diagonal
trivialization.  This proves \eqref{foundations:evaluation}.

\emph{Return to the full bundle stack.}
Transport compact generation and duality across
\eqref{foundations:cutoff-equivalence}.  The spectator version of that
equivalence identifies products and diagonal kernels, so the same
formulas apply.  The renormalized pairing may equivalently be computed
on the common cutoff open; it is independent of enlarging that open.
For local coherence, one need not assert that an arbitrary open star
image preserves coherence.  Given a finite-type smooth chart of $Y_H$,
choose a larger quasi-compact HN open $V$ containing both its image and
$U$.  The same cutoff argument gives $D_\tau(V)\simeq D_\tau(U)$.
A compact object is thus compact on $V$, where the preceding Spencer
argument proves its coherence on the chosen chart.  This proves local
coherence on the full stack.
\end{proof}

In particular, $A'$ and $B'$ realize the ordinary categorical duals of
$A$ and $B$.  Formula~\eqref{foundations:evaluation} still uses the
star-diagonal kernel.  The shriek-diagonal kernel
$\Delta_!k_Y$, which defines pseudo-identity, is a separate object;
proving that it gives an equivalence is the next step.
